\reportsection{06}{Terminal}{node-pty: Compiling What Windows Would Not Ship}

\unithead{Mechanism}

\reppar{This unit is \textbf{build-time provisioning, not a bundle sed patch}: it touches no minified JavaScript. Its job on main is to put a Linux-native \code{pty.node} where the Windows-installer \code{app.asar} expects one, so the Code tab's terminal (\code{startShellPty}, which dynamically imports \code{node-pty} as an optional dependency) can spawn shells. Four pieces cooperate: \code{install\_node\_pty()} in \code{cowork.sh} either consumes a pre-built tree via \code{--node-pty-dir} or npm-installs into an isolated directory; \code{patch\_app\_asar()} injects \code{optionalDependencies['node-pty']} so the import resolves; \code{finalize\_app\_asar()} packs with \code{--unpack '**/*.node'}; and \code{nix/node-pty.nix} builds node-pty v1.1.0 from source with three commented upstream workarounds.}

\unithead{Origin}

\reppar{The origin is \#152 (add node-pty for terminal integration), merged in January 2026: macOS shipped node-pty as an optional dependency, but the repackaged Windows \code{app.asar} carried no Linux \code{pty.node}. No motivating issue exists; the PR reads as self-initiated, though that is inference. @typedrat's Nix flake PR \#266 (March 2026) briefly consumed a nixpkgs \code{node-pty} that turned out not to exist, replaced within the PR by \code{nix/node-pty.nix} and the \code{--node-pty-dir} wiring.}

\begin{chronology}
2026-01 & \#152, \code{3591392} & Origin: \code{install\_node\_pty()} and the \code{optionalDependencies['node-pty']} edit land directly in \code{build.sh}; motivation is the macOS optional-dependency precedent, since the repackaged Windows \code{app.asar} shipped no Linux \code{pty.node}; no motivating issue number, reads as self-initiated. \\
2026-03 & \#266, \code{caa58ca} & Nix flake integration (@typedrat): a nixpkgs \code{node-pty} input is added, then removed within the same PR after proving not to exist; \code{nix/node-pty.nix} built from source instead (v1.1.0, three commented workarounds), wired through \code{--node-pty-dir}. \\
2026-04 & \#421, \code{3150477} & Asar-manifest unpack fix (@Joost-Maker): no manifest entry for \code{node-pty/build/} meant Electron's \code{.unpacked} redirect never fired, and the require died with \code{MODULE\_NOT\_FOUND} on every Code-tab shell session; \code{install\_node\_pty()} now stages \code{build/}, and \code{finalize\_app\_asar()} packs with \code{--unpack '**/*.node'}. \\
2026-04 & \#432/\#438, \code{50b10ed}, \code{e92aea1} & Nix permissions (@typedrat): the \#421 \code{--unpack} pass preserved Nix-store read-only bits, breaking the follow-up copy; \#432 chmods the extracted files, \#438 retires the chmod for \code{cp --no-preserve=mode} at the source. \\
2026-05 & \#401, \#598/\#597 & Fedora ships Windows node-pty binaries (\#401): the origin's soft-warning design let a failed npm install through, shipping \code{winpty.dll} and PE32+ binaries unchanged; @JoshuaVlantis closes both halves via \#598 (fail loudly on install failure) and \#597 (wipe the upstream Windows tree before staging). \\
2026-06 & \#727/\#728, \#761 & Open (@EtherAura): node-pty 1.1.0 also forks the shell through a separate \code{build/Release/spawn-helper} executable the legacy build never shipped (\#727), and the bundle hardcodes \code{powershell.exe} (\#728). @LiukScot's \#761 targets \code{nix/node-pty.nix} and \code{cowork.sh}, files the rebase deletes or parks; its disposition is unrecorded. \\
\end{chronology}

\methodbanner{\textbf{Fate under the rebase.} \bdel{} The matrix verdict is an unconditional "delete": "Prebuilt \code{prebuilds/linux-x64/pty.node} ships in \code{app.asar.unpacked}." The official 1.17377.2 \code{.deb} carries a prebuilt Linux ELF, in node-pty's newer \code{prebuilds/} layout rather than the legacy \code{build/Release/} path, verified by \code{probe\_node\_pty()} in the audit tool. Commit d9cef9e deletes the whole unit: \code{nix/node-pty.nix} gone, \code{install\_node\_pty} nowhere, \code{--node-pty-dir} dropped. The new repack derives its \code{--unpack} glob from the shipped \code{app.asar.unpacked} tree with a post-pack set-equality check, preserving the official placement. One hedge: the delete rests on byte presence of the ELF, not a live shell-spawn test, so whether the official build moots \#727, \#728 is unverified.\\[3pt]\textbf{Affects.} All Linux; build-time provisioning, DE-agnostic, format-general (deb, RPM, Nix).}
